% !TEX program = pdflatex
\documentclass[11pt]{article}

\usepackage[margin=1in]{geometry}
\usepackage{amsmath,amssymb,amsthm}
\usepackage{hyperref}
\usepackage{enumitem}

\title{Symmetry Constraints and Saturation in the Hypercube}
\author{Amos Jay Maley}
\date{\today}

% ---------- Theorem environments ----------
\theoremstyle{plain}
\newtheorem{theorem}{Theorem}[section]
\newtheorem{lemma}[theorem]{Lemma}
\newtheorem{corollary}[theorem]{Corollary}

\theoremstyle{definition}
\newtheorem{definition}[theorem]{Definition}
\newtheorem{remark}[theorem]{Remark}

% ---------- Notation ----------
\newcommand{\Qn}{Q_n}
\newcommand{\Aut}{\mathrm{Aut}}

\begin{document}
\maketitle

\begin{abstract}
We study square-free saturated subgraphs of the $n$-dimensional hypercube under symmetry constraints.  
Introducing a symmetry-restricted saturation parameter, we show that full automorphism invariance precludes nontrivial saturation for $n\ge3$.  
We then identify the largest vertex-transitive symmetry group compatible with saturation, determine the exact minimum edge count under this symmetry, and classify all extremal configurations.  
This establishes a sharp boundary between symmetry and maximal square-avoidance in the hypercube.
\end{abstract}

\section{Introduction}

Let $\Qn$ denote the $n$-dimensional hypercube.  
A copy of $Q_2$ in $\Qn$ is a $4$-cycle induced by two coordinates.  
A subgraph $G\subseteq \Qn$ is \emph{$(Q_n,Q_2)$-saturated} if it contains no $4$-cycle, but the addition of any missing cube-edge creates one.

Saturated subgraphs of the hypercube have been studied extensively without symmetry restrictions.  
In this paper we initiate a complementary viewpoint: \emph{extremal saturation under symmetry constraints}.

\subsection{A symmetry-restricted saturation parameter}

Let $H\le \Aut(\Qn)$ be a subgroup of the hypercube automorphism group.

\begin{definition}
Define
\[
\mathrm{sat}_H(\Qn,Q_2)
\;=\;
\min\{|E(G)|:\;
G\subseteq \Qn,\;
G\text{ is $(Q_n,Q_2)$-saturated and $H$-invariant}\}.
\]
\end{definition}

The goal is to understand how increasing symmetry restricts maximal square-avoidance.

\section{Full Symmetry Obstruction}

We begin with a necessary boundary observation.

\begin{lemma}[Full symmetry obstruction]
For $n\ge3$, no proper $\Aut(\Qn)$-invariant $(Q_n,Q_2)$-saturated subgraph of $\Qn$ exists.
\end{lemma}

\begin{proof}
The automorphism group $\Aut(\Qn)$ acts transitively on edges.  
Hence any $\Aut(\Qn)$-invariant edge set is either empty or all of $E(\Qn)$.  
A saturated subgraph must be nonempty and proper, which is impossible under full invariance.
\end{proof}

\begin{remark}
This lemma is not presented as a deep result, but as a structural obstruction motivating the search for the largest symmetry compatible with saturation.
\end{remark}

\section{Maximal Symmetry Compatible with Saturation}

Let $S_n$ denote the group acting on $\Qn$ by permuting coordinates (without bit-flips).

\begin{lemma}
Any vertex-transitive subgroup $H\le\Aut(\Qn)$ strictly containing $S_n$ is edge-transitive.
\end{lemma}

\begin{proof}
Any element of $\Aut(\Qn)$ outside $S_n$ introduces a nontrivial bit-flip.  
Together with coordinate permutations, this generates the full automorphism group, which is edge-transitive.
\end{proof}

\begin{corollary}
Among vertex-transitive subgroups of $\Aut(\Qn)$, $S_n$ is maximal with more than one edge orbit.
\end{corollary}

Thus $S_n$ is the largest vertex-transitive symmetry group under which nontrivial invariant edge-sets exist.

\section{$S_n$-Invariant Saturated Subgraphs}

\subsection{Orbit structure}

Under the action of $S_n$, edges of $\Qn$ decompose into $n$ orbits, one for each coordinate direction.  
Any $S_n$-invariant subgraph is determined by the subset of coordinate orbits it contains.

\subsection{Square forcing}

Each $4$-cycle in $\Qn$ corresponds to a pair of coordinates.  
Saturation imposes explicit constraints: omitting too many coordinate directions leaves a square unguarded.

\section{Classification of Extremal Configurations}

\begin{theorem}[Exact minimizers under $S_n$]
For $n\ge3$:
\begin{enumerate}[label=(\roman*)]
\item $\mathrm{sat}_{S_n}(\Qn,Q_2)$ is attained.
\item The minimum edge count equals
\[
\mathrm{sat}_{S_n}(\Qn,Q_2)
=
(n-1)2^{n-1}.
\]
\item Every minimizer is obtained by omitting exactly one coordinate direction (up to symmetry), with a fixed parity orientation.
\item No other $S_n$-invariant saturated subgraph attains this minimum.
\end{enumerate}
\end{theorem}

\begin{proof}
Omitting two or more coordinate directions leaves a square with both directions absent, violating saturation.  
Omitting exactly one direction yields a square-free subgraph in which every missing edge is witnessed by a $4$-cycle.  
Counting edges gives the stated minimum.  
Uniqueness follows from orbit considerations.
\end{proof}

\section{Cost of Symmetry}

Let $\mathrm{sat}(\Qn,Q_2)$ denote the unrestricted saturation number.

\begin{corollary}
For $n\ge3$,
\[
\mathrm{sat}_{S_n}(\Qn,Q_2)
>
\mathrm{sat}(\Qn,Q_2),
\qquad
\text{but}\qquad
\mathrm{sat}_{S_n}(\Qn,Q_2)
=
\Theta(\mathrm{sat}(\Qn,Q_2)).
\]
\end{corollary}

\begin{remark}
Thus symmetry imposes a genuine extremal cost, while preserving the correct order of magnitude.
\end{remark}

\section{Boundary Sharpness}

\begin{theorem}
Let $H\le\Aut(\Qn)$ be vertex-transitive.
If $H\not\subseteq S_n$, then either:
\begin{enumerate}[label=(\roman*)]
\item $\mathrm{sat}_H(\Qn,Q_2)$ is undefined, or
\item every $H$-invariant minimizer coincides with an $S_n$-orbit extremizer.
\end{enumerate}
\end{theorem}

\section{Small Dimensions}

The cases $n=1,2$ are degenerate.  
For $n=3$, all claims can be verified by direct inspection.

\section{Conclusion}

We identify the largest vertex-transitive symmetry compatible with maximal square-avoidance in the hypercube and classify all extremal configurations under that symmetry.  
This introduces symmetry-restricted saturation as a new extremal lens for problems on structured graphs.

\end{document}
